% !TeX root = ../../Book1.tex
% 纲要标题：### 6.1 曲面与局部组合结构
% 纲要位置：几何学计划纲要.md，第 319 行。
% 正文状态：初稿；尚未正式审核。

% 纲要内容开始
% 
% * 内点、边界点与连通分支
% * 局部圆盘的相容缩小
% * 定向的局部定义
% 
% 纲要内容结束

\section{曲面与局部组合结构}
\label{g1:sec:0601}

球面、环面和圆盘在大部分点附近看起来都像一小块平面，差别在于这些小块怎样接成整体。圆盘的边缘还需要另一种局部模型：只能向一侧延伸的半平面。本节先把这两种点的区别、局部邻域的控制和定向说清楚。这里没有切向量，也没有坐标变换的导数；判断方向将使用上一章已经建立的绕数。

\subsection{内点、边界点与连通分支}
\label{g1:subsec:0601:01}

\begin{definition}[拓扑曲面]\label{g1:def:topological-surface}
	设 \(M\) 为 Hausdorff、第二可数的拓扑空间。若每个点都有一个开邻域 \(U\)，同胚于闭半平面
	\[
		\mathbb H^2=\{(x,y)\in\mathbb R^2:y\ge0\}
	\]
	中的相对开集，则称 \(M\) 为允许带边界的\term[tuopuqumian]{拓扑曲面}[topological surface]。这样的同胚 \(\varphi:U\to\varphi(U)\subseteq\mathbb H^2\) 称为\term[zuobiaotu]{坐标图}[coordinate chart]。当 \(\varphi(p)\) 的第二坐标大于零时，称 \(p\) 为内点；等于零时，称 \(p\) 为曲面的边界点。所有边界点组成的集合记为 \(\partial M\)，内点集记为 \(M^\circ\)。
\end{definition}

定义暂时用某张图区分两类点。必须证明它不随换图改变，否则 \(\partial M\) 就不是空间本身的性质。平面圆盘去掉中心后有不能缩掉的环路，半圆盘去掉直边中心后却没有；这提供了一个不依赖坐标名称的判断。

\begin{proposition}[边界的内在性]\label{g1:prop:boundary-invariance}
	在允许带边界的拓扑曲面中，一个点不可能在一张坐标图中为内点、在另一张图中为边界点。任一曲面同胚 \(f:M\to N\) 因而满足 \(f(\partial M)=\partial N\)。内点集开，边界集闭；\(\partial M\) 在子空间拓扑下是一维无边界拓扑流形。
\end{proposition}
\begin{proof}
	假设 \(p\) 同时有两种图。先在内点图中取以 \(p\) 为中心的坐标开圆盘 \(B\)，使它包含于两图的共同定义域；再在边界图中取以 \(p\) 为直边中心的坐标开半圆盘 \(D\subset B\)；最后取内点图中更小的同心圆盘 \(B'\subset D\)。选 \(z\in B'\setminus\{p\}\) 为共同基点。包含映射给出
	\[
		\pi_1(B'\setminus\{p\},z)\longrightarrow
		\pi_1(D\setminus\{p\},z)\longrightarrow
		\pi_1(B\setminus\{p\},z).
	\]
	两端的群都是 \(\mathbb Z\)，而复合为同构：同心穿孔圆盘沿半径收缩到小圆周，包含映射把绕一周送到绕一周，所用计算是\cref{g1:thm:circle-fundamental}。中间的群却平凡，因为穿孔半圆盘沿半径形变收缩到一条闭半圆弧，后者同胚于闭区间而可缩。这与复合同态为同构矛盾。

	同胚将坐标图带到坐标图，故保持两种点。内点有完全由内点组成的小圆盘邻域，所以 \(M^\circ\) 开，\(\partial M\) 闭。在边界图中，刚证的内在性说明边界恰对应直边上的一个开区间；这些限制图给出 \(\partial M\) 的一维无边界流形结构。
\end{proof}

这里的 \(\partial M\) 是曲面的内在边界，不是任意嵌入下的子集边界。例如圆盘 \(D^2\) 的曲面边界是 \(S^1\)，球面 \(S^2\) 的曲面边界为空；但把 \(S^2\) 放在 \(\mathbb R^3\) 中，它作为子集的拓扑边界是整个 \(S^2\)。第四章的 \(\partial_X A\) 仍指后一种依赖环境的边界。紧且边界为空的曲面称为\term[biqumian]{闭曲面}[closed surface]；“闭”在这里同时包含紧性，不能只理解为某个环境中的闭子集。

\begin{proposition}[分支与边界圆周]\label{g1:prop:boundary-circles}
	允许带边界的拓扑曲面局部道路连通，其连通分支就是道路分支，并且开。紧曲面只有有限个连通分支，其边界是有限多个彼此不交的圆周的并。
\end{proposition}
\begin{proof}
	小圆盘与小半圆盘都道路连通，且可以任意缩小，故曲面局部道路连通。由\cref{g1:thm:locally-path-components}，两种分支相同并且开。若曲面紧，用这些开分支覆盖它，只能有有限多个分支。边界闭，因而在紧曲面中紧；上命题又说明它是一维无边界流形。

	下面补上所需的一维事实，不预设一维流形分类。设 \(C\) 为紧的一维无边界流形。在每一点选取一段嵌入的闭区间，使该点落在其中部；由紧性，有限条闭区间的内部覆盖 \(C\)。把这些区间的所有端点组成有限集 \(P\)。每条区间去掉 \(P\) 后只有有限个开小段。因此 \(C\setminus P\) 只有有限个连通分支：它由有限个这样的连通小段覆盖。

	任取一个分支 \(J\) 和其中一点，选择覆盖该点的一个区间内部 \(U\)。由于 \(U\) 的两个端点都已删去，\(U\setminus P\) 在 \(C\setminus P\) 中既开又闭，故 \(J\) 完全落在其中一段内。于是 \(J\) 本身就是一个开区间，它的两端在 \(P\) 中。这样 \(C\) 成为有限图；每个顶点的区间邻域恰有左右两个半边，所以顶点的度数为二。从任一边出发沿唯一的下一条边继续，有限性迫使回到起点，且不能遗下同一连通分支中的其他边。因此每个连通分支是一个有限循环图，同胚于 \(S^1\)。应用于 \(C=\partial M\) 即得结论；空边界对应零个圆周。
\end{proof}

这个结论允许我们先分类连通紧曲面，再取有限不交并。以后用 \(b\) 记边界圆周数；它是曲面同胚必须保留的非负整数。

\subsection{局部圆盘的相容缩小}
\label{g1:subsec:0601:02}

一张坐标图不能覆盖整个环面，但在已有坐标图里，可以预留一圈余地后再动手修改曲线。这种“较小集合的闭包仍在较大坐标图中”的控制，会防止局部切割碰到图的边缘。此时相容指拓扑包含和覆盖关系，不预设坐标变换可微。

\begin{lemma}[带余量的坐标圆盘]\label{g1:lem:nested-surface-charts}
	设 \(M\) 为允许带边界的拓扑曲面，\(p\in M\)，\(O\) 为 \(p\) 的开邻域。可以在同一坐标图中选取三个相似的闭圆盘或闭半圆盘 \(K_1,K_2,K_3\)，使
	\[
		p\in\operatorname{Int}_M K_1,\qquad
		K_1\subset\operatorname{Int}_M K_2,\qquad
		K_2\subset\operatorname{Int}_M K_3,\qquad K_3\subset O.
	\]
	半圆盘的直边位于 \(\partial M\)。若 \(M\) 紧，可以选有限组这样的集合，使最小各片的相对内部覆盖 \(M\)。
\end{lemma}
\begin{proof}
	先限制坐标图到 \(O\) 内。在内点处，将坐标中心平移到原点，选闭球完全落在坐标像内，再取三个严格递增的小半径。在边界点处，对半平面取以直边上的坐标点为中心的三个闭半球。其直边中部在 \(M\) 中仍是相对内部点，因为半球本来就只需要包含半平面内的邻域。

	这些闭圆盘或半圆盘的坐标原像紧；\(M\) 为 Hausdorff，所以它们在 \(M\) 中闭，闭包关系与坐标中的关系一致。这说明每一级都保留了所需余量。对所有 \(p\) 先作选择，最小各片的内部成为开覆盖，紧性再给出有限子覆盖。
\end{proof}

每一点都有紧邻域，所以曲面局部紧。由\cref{g1:cor:lch-second-countable}，Hausdorff、第二可数与局部紧性保证存在相容距离；本章在需要控制网格大小时固定其中任意一个。距离的数值不是新的曲面不变量，它只帮助把“充分小”写成统一的有限条件。特别是对紧曲面的有限坐标覆盖，\cref{g1:lem:lebesgue-number}提供正数 \(\delta\)，使直径小于 \(\delta\) 的集合完全落入某一坐标图。

\ProseParagraphBreak

还需注意：两张图中的圆看起来都很规则，它们在重叠处却可能无限次相交。坐标变换目前只是同胚，不能用“稍微移动就横截”代替证明。下一节将在平面中建立简单曲线和简单弧的具体控制，再把有限圆盘覆盖整理成有限曲线网。有限覆盖本身尚不等于有限三角剖分。

\subsection{定向的局部定义}
\label{g1:subsec:0601:03}

在欧氏平面中，“逆时针”有确定意义。换到另一张拓扑坐标图后，一个小圆可能变成非常不规则的简单闭曲线，但它围绕像点的绕数仍有意义。对不经过 \(a\in\mathbb R^2\) 的环路 \(\gamma\)，定义
\[
	w(\gamma,a)=w\left(t\longmapsto
	\frac{\gamma(t)-a}{\lVert\gamma(t)-a\rVert}\right).
\]
右边是\cref{g1:thm:circle-fundamental}的圆周绕数。起点不必为 \(1\)：把角度提升的终值减去初值即可，改变初始角度不会改变这个整数。移动基点的环路同伦仍保持总角度增量，亦可由\cref{g1:prop:moving-basepoint}和 \(\mathbb Z\) 的交换性得到。

\begin{proposition}[平面换图的局部方向符号]\label{g1:prop:orientation-local-sign}
	设 \(U,V\subseteq\mathbb R^2\) 为开集，\(f:U\to V\) 为同胚。对 \(x\in U\)，取以 \(x\) 为中心、闭包包含于 \(U\) 的小圆，按逆时针参数化为 \(\gamma_x\)。整数
	\[
		\varepsilon_f(x)=w(f\circ\gamma_x,f(x))
	\]
	与小圆半径无关，只能等于 \(1\) 或 \(-1\)，且关于 \(x\) 局部常值。对可复合的同胚，有
	\[
		\varepsilon_{g\circ f}(x)=\varepsilon_g(f(x))\,\varepsilon_f(x),
		\qquad \varepsilon_{\operatorname{id}}=1.
	\]
\end{proposition}
\begin{proof}
	同心小圆可沿半径在 \(U\setminus\{x\}\) 中变形到彼此，因此它们的像有相同绕数。为证明这个数的绝对值为一，取以 \(x\) 为中心的开圆盘 \(B\)，再取以 \(f(x)\) 为中心的圆盘 \(D\)，使 \(\overline D\subset f(B)\)，最后取以 \(x\) 为中心的圆盘 \(B'\)，使 \(f(B')\subset D\)。穿孔空间上的同态组成
	\[
		\pi_1(B'\setminus\{x\})\xrightarrow{\,f_*\,}
		\pi_1(D\setminus\{f(x)\})\longrightarrow
		\pi_1(f(B)\setminus\{f(x)\}).
	\]
	这里各群都是无限循环群：前两者由径向收缩得到，第三者由同胚 \(f\) 与 \(B\setminus\{x\}\) 得到。复合同态是同构，因为它就是 \(f\) 复合同心穿孔圆盘的包含同构。把两个同态写成整数乘法，其乘数的乘积为 \(1\) 或 \(-1\)，所以两个乘数都为 \(1\) 或 \(-1\)。第一个乘数正是所定义的绕数。基点可沿相应小区域中的道路对齐；循环群的内自同构平凡，不影响这些整数。

	证明局部常值时，固定一条围住 \(x\) 的坐标圆 \(C\)，其闭圆盘包含于 \(U\)。对圆盘内部任何 \(y\)，\(C\) 都代表绕 \(y\) 一周的类：以 \(y\) 为中心向一个充分小的圆作径向收缩，所有径向线段留在这个凸圆盘中。因此
	\[
		\varepsilon_f(y)=w(f\circ C,f(y)).
	\]
	当 \(y\) 沿圆盘内部的道路移动时，\(f(y)\) 始终避开 \(f(C)\)，上式的归一化方向连续变化，给出圆周环路的同伦。于是整数不变，证明局部常值。

	最后在 \(x,f(x),g(f(x))\) 处逐级取充分小的圆盘。穿孔小邻域中的基本群均由正向小圆指定生成元；\(f_*\) 和 \(g_*\) 分别是乘以相应符号，复合便乘以两者之积。恒等映射显然把正向圆保留为正向圆。
\end{proof}

若换图涉及边界，先在半圆盘的内点部分使用上述符号。这个部分连通，所以符号恒定；把相同符号赋给其直边点。缩小半圆盘不会改变符号，乘法公式也通过内点的公式延续到边界。这里不把半圆盘边界上的点当作平面内点来计算绕数。

\begin{definition}[局部定向与曲面的定向]\label{g1:def:surface-orientation}
	在固定点 \(p\)，把一张坐标图 \(\varphi\) 和一个符号 \(s_\varphi\in\{1,-1\}\) 配在一起；若另一对 \((\psi,s_\psi)\) 满足 \(s_\psi=\varepsilon_{\psi\circ\varphi^{-1}}(\varphi(p))s_\varphi\)，便把这两对视为等价。符号的乘法公式保证这是等价关系，每个点恰有两个等价类，称为该点的两个局部定向。它们记录同一方向在不同坐标中的写法。
	设 \(M\) 为允许带边界的拓扑曲面。在每个点的坐标图中标记一个符号 \(s_\varphi\in\{1,-1\}\)，其中符号在图内局部常值；若在任意两图交叠处都满足
	\[
		s_\psi(p)=\varepsilon_{\psi\circ\varphi^{-1}}(\varphi(p))\,s_\varphi(p),
	\]
	并将满足此关系的局部描述视为同一选择，则称这些相容选择为 \(M\) 的一个\term[dingxiang]{定向}[orientation]。存在定向的曲面称为\term[kedingxiang]{可定向}[orientable]，否则称为不可定向。直观上，符号 \(1\) 选定图中的逆时针方向，符号 \(-1\) 选定相反方向。
\end{definition}

只需在覆盖 \(M\) 的一组坐标图上给出相容符号：乘法公式保证它们唯一地确定所有其他图中的符号。这个定义不要求选择一个全局坐标系；它要求的是各处的局部方向能接在一起。对连通可定向曲面，定向恰有两个，因为两个定向在每点的符号比值为局部常值的 \(\pm1\)，在连通空间上只能恒定。

\begin{example}[平面方向与环面的定向]\label{g1:exm:torus-local-orientation}
	平移与旋转把正向小圆送到正向小圆，故方向符号为 \(1\)；反射 \((x,y)\mapsto(x,-y)\) 的符号为 \(-1\)。在环面模型 \(\mathbb R^2/\mathbb Z^2\) 中，充分小的投影圆盘作为坐标图，重叠处的坐标变换为整数平移。因此所有这些图均选符号 \(1\) 就得到环面的定向。将全部符号同时取反，得到另一个定向。
\end{example}

对平面中带边界的区域，正定向规定沿边界行走时区域留在左侧。若取半平面 \(y\ge0\) 的正定向，直边上的正向就是 \(x\) 增大的方向。定向相容的换图保留这个选择：把一段直边与内侧一条简单弧合成小半圆盘的边界，边界的行走方向由内点的正向小环路决定，换图符号为正时这个方向不变。下一节的简单弧直化会给出这个边界描述的圆盘坐标。反转曲面定向时，各边界圆周的诱导方向也同时反转。
